% ID: FIS-ROT-ANG-001
% Titolo: Sei ragazze su una piattaforma rotante
% Disciplina: Fisica
% Area: Dinamica rotazionale
% Argomento: Conservazione del momento angolare
% Sottoargomento: Variazione del momento d'inerzia in un sistema ridotto
% Classe: Quarta liceo scientifico
% Difficolta: 3
% Tipo: problema_numerico
% Risultato: due ragazze; omega_f circa 2,33 rad/s
% Tempo_stimato: 10 min
% Competenze: modellizzazione, conservazione del momento angolare, momento d'inerzia
% Prerequisiti: moto circolare, momento d'inerzia, velocita angolare
% Tag: fisica, dinamica_rotazionale, momento_angolare, momento_inerzia, sistema_ridotto
% Fonte: esercizio_proposto_utente
% Autore: Riccardo Carli
% Licenza: CC-BY-SA-4.0
% Simulazione: rotational_platform

\begin{esercizio}
Una piattaforma circolare di diametro \(2{,}00\,\mathrm{m}\) e massa
\(200\,\mathrm{kg}\) ruota senza attrito attorno al proprio asse con velocita
angolare iniziale
\[
\omega_i=1{,}80\,\mathrm{rad/s}.
\]
Sei ragazze di massa \(50{,}0\,\mathrm{kg}\) ciascuna si trovano, equidistanti,
alla distanza \(0{,}850\,\mathrm{m}\) dall'asse.

Nel modello didattico dell'esercizio si considera il momento angolare del
sistema ridotto formato dalla piattaforma e dalle ragazze che vi restano come
costante. La piattaforma e modellizzata come un disco pieno e ogni ragazza come
una massa puntiforme. Quante ragazze devono lasciare la piattaforma affinche la
velocita angolare diventi circa \(2{,}34\,\mathrm{rad/s}\)?

Nota sul modello: una persona che lascia realmente la piattaforma puo
portare con se momento angolare. L'ipotesi di conservazione applicata qui al
sistema ridotto e una semplificazione didattica, non una descrizione completa
della fase di salto.
\end{esercizio}

\begin{soluzione}
Il raggio della piattaforma e \(R=1{,}00\,\mathrm{m}\). Il suo momento
d'inerzia vale
\[
I_\text{piattaforma}=\frac{1}{2}MR^2
=\frac{1}{2}\cdot 200\cdot(1{,}00)^2
=100\,\mathrm{kg\,m^2}.
\]

Ogni ragazza contribuisce con
\[
I_\text{ragazza}=mr^2
=50{,}0\cdot(0{,}850)^2
=36{,}125\,\mathrm{kg\,m^2}.
\]

Con sei ragazze il momento d'inerzia iniziale e
\[
I_i=100+6\cdot36{,}125
=316{,}75\,\mathrm{kg\,m^2},
\]
e il momento angolare di riferimento assunto dal modello e
\[
L_\text{rif}=I_i\omega_i
=316{,}75\cdot1{,}80
=570{,}15\,\mathrm{kg\,m^2/s}.
\]

Se lasciano la piattaforma \(k\) ragazze, ne restano \(6-k\) e
\[
\omega_f(k)=\frac{L_\text{rif}}
{100+(6-k)\cdot36{,}125}.
\]

Con due ragazze in meno ne restano quattro:
\[
I_f=100+4\cdot36{,}125
=244{,}5\,\mathrm{kg\,m^2},
\]
\[
\omega_f=\frac{570{,}15}{244{,}5}
\simeq2{,}332\,\mathrm{rad/s}.
\]
Il valore e circa \(2{,}33\)--\(2{,}34\,\mathrm{rad/s}\), secondo gli
arrotondamenti. Devono quindi lasciare la piattaforma \textbf{due ragazze},
nel modello didattico dichiarato dal testo.
\end{soluzione}
